\documentclass{article}
\usepackage{graphicx}     % images
\usepackage{amsmath}      % equations
\usepackage{booktabs}     % nice tables
\usepackage{array}
\usepackage{xcolor}
\usepackage[hidelinks]{hyperref}

% --- placeholder box for figures to be added later ---
\newcommand{\plotph}[1]{%
  \begin{center}\fbox{\parbox[c][3.5cm][c]{0.85\linewidth}{%
  \centering\itshape [Plot placeholder: #1]}}\end{center}}

\title{Zebrafish bout discretization and hidden-Markov modelling from the
       semi-closed-loop visual-stimulus assay}
\author{}
\date{}

\begin{document}
\maketitle

\section{Data}

\subsection{Experimental paradigm}
The dataset (\texttt{GsizeMotionSingle}) comes from a semi-closed-loop
behavioural assay (Perrino PhD thesis, Sect.~4.4): juvenile zebrafish
(15--21~dpf) swim freely, each isolated in a 10~cm-diameter circular Petri
dish, arranged in a $3\times5$ grid so that \textbf{15 fish are recorded
simultaneously}. Behaviour is filmed from above with an infrared camera at
\textbf{30~frames\,s$^{-1}$}; the full session lasts $\sim$180~min.

The visual stimuli are groups of black dots (2~mm diameter) projected from
below. They are designed to \emph{simulate conspecifics}: at stimulus onset the
dot group is placed left or right of the fish (relative to the fish centroid),
then follows a predefined open-loop trajectory within a $2.5\times2.5$~cm area.
Stimuli vary in:
\begin{itemize}
  \item \textbf{group size}: 1, 2, 4, 6 or 8 dots, presented on a single side
        (hence ``Single'');
  \item \textbf{speed}: 0~mm\,s$^{-1}$ (static) or 3.5~mm\,s$^{-1}$ (moving,
        bout-like).
\end{itemize}
Each trial is a 60~s stimulus presentation followed by a 50~s pause; conditions
are repeated 6 times per side in randomised blocks.

\plotph{Example egocentric trajectory of one fish before/after stimulus onset
        with the projected dot group (cf.\ notebook diagnostics).}

\subsection{Raw acquisition files}
Three raw files are produced per session (camera origin at the top-left;
distances in pixels unless stated):

\begin{table}[h]\centering\small
\begin{tabular}{@{}lp{0.62\linewidth}@{}}
\toprule
\textbf{File} & \textbf{Content}\\
\midrule
\texttt{PositionTxt} & One row per frame. First 45 columns = $(x,y,\text{orientation})$
  for each of the 15 fish (pixels; orientation in radians, measured
  \emph{clockwise}). Remaining columns = stimulus dot data.\\
\texttt{Trajectory}  & One row per frame. Dot coordinates $(x,y)$ in
  \emph{millimetres relative to the fish centroid}, plus stimulus name, size and
  colours. Convention: left $\to x>0$, right $\to x<0$.\\
\texttt{ROIdef}      & One row per dish: \texttt{xoff, yoff, diameter, x, y,
  radius} (dish geometry in pixels).\\
\bottomrule
\end{tabular}
\end{table}

\paragraph{Stimulus naming.} e.g.\ \texttt{CLsemi1800\_g4-0\_s3.5-3.5}:
\texttt{CLsemi} = semi-closed loop; \texttt{1800} = presentation length in frames
($1800/30 = 60$~s); \texttt{g4-0} = 4 dots left, 0 right; \texttt{s3.5-3.5} =
speed (mm\,s$^{-1}$). When no stimulus is shown the field is \texttt{none}
(pause).

\subsection{Preprocessing pipeline}
The raw files are transformed into analysis-ready per-fish dataframes by
\texttt{src/transformed\_data\_pipeline.py}
(\texttt{generate\_transformed\_datasets}). The main steps are:
\begin{enumerate}
  \item \textbf{Naming \& sanity check}: assign column names
        (\texttt{f\{N\}\_x/\_y/\_ori}, \texttt{d\{i\}\_x/\_y/\_stim}) and verify
        that the stimulus labels agree between \texttt{PositionTxt} and
        \texttt{Trajectory}.
  \item \textbf{Trial structure}: derive trial / stimulus-block indices; tag each
        frame's \texttt{condition} (a \texttt{CLsemi\dots} label, or
        \texttt{inter\_stim\_pause}); separate out the end-of-session
        single-dot ``knot'' positive-control (\texttt{bout}).
  \item \textbf{Stimulus geometry}: convert dot coordinates from
        fish-centred mm to absolute screen coordinates using the fish position
        at block onset; count dots left/right of the fish.
  \item \textbf{Coordinate normalisation} (\texttt{transform\_coordinates}):
        invert the $y$ axis ($y \mapsto \text{diameter}-y$) and convert
        orientation to counter-clockwise ($\theta \mapsto 2\pi-\theta$), so the
        working frame is CCW with $0 =$ right.
  \item \textbf{Assembly}: build \texttt{abs\_dict =
        \{"f0":df0,\dots,"f14":df14\}}, one dataframe per fish. A polar
        representation (\texttt{d\{k\}\_angle}, \texttt{d\{k\}\_radius}) is also
        produced.
\end{enumerate}
The result is cached as the per-fish CSVs
\texttt{*\_f\{N\}\_transformed\_abs.csv} (equivalently \texttt{abs\_dict.pkl}).
Stacking these with \texttt{functions.stack\_dict\_pipeline} yields the long
dataframe \texttt{df\_abs\_raw} (columns \texttt{f\{N\}\_*} renamed to
\texttt{f\_*}, plus a \texttt{fish\_id} column).

\subsection{Working dataframe}
We segment bouts directly from the per-fish dataframes (one file $=$ one fish,
$f0\dots f14$). Each has one row per frame with the columns:

\begin{table}[h]\centering\small
\begin{tabular}{@{}lll@{}}
\toprule
\textbf{Column} & \textbf{Meaning} & \textbf{Units / type}\\
\midrule
\texttt{frame}            & frame index (30~fps)                 & integer\\
\texttt{stimulus}         & condition active this frame          & label\\
\texttt{f\{N\}\_x},\,\texttt{f\{N\}\_y} & fish centroid              & \textbf{pixels} (integer)\\
\texttt{f\{N\}\_ori}      & body orientation (CCW, $0=$ right)   & radians $[0,2\pi)$\\
\texttt{d\{i\}\_x},\,\texttt{d\{i\}\_y} & dot $i$ position ($i=0..7$) & pixels\\
\texttt{d\{i\}\_stim}     & dot $i$ stimulus label               & label\\
\bottomrule
\end{tabular}
\end{table}

\paragraph{Important conventions.}
\begin{itemize}
  \item \textbf{Positions are integer pixels} -- sub-pixel precision was lost
        upstream, which quantises short bouts (mitigated later by smoothing
        \emph{for detection only}).
  \item \textbf{The pixel$\to$mm scale is per dish.} The pipeline constant
        \texttt{MM\_TO\_PX}$=3.67$ is a single \emph{nominal camera} scale (the
        repo's own comment: a 50~mm-radius dish ``should'' be
        $50\times3.67\approx183.5$~px), and it is applied \emph{only to the
        stimulus dots}, never to the fish positions. Because each dish is imaged
        at a different magnification, the real camera scale is recovered per fish
        from
        \[ s_f \;=\; \frac{50~\text{mm}}{\texttt{radius}_f~[\text{px}]}, \]
        where $\texttt{radius}_f$ is that dish's radius in \texttt{ROIdef}
        (physical dish radius $=50$~mm). The 15 dishes span
        $2.5$--$4.7$~px\,mm$^{-1}$ (mean $\approx 3.33$, not $3.67$), so a single
        global value would mis-scale a given fish's distances by up to
        $\sim\pm45\%$. Only \emph{distances} depend on $s_f$; angles
        ($\delta\theta$) do not.
\end{itemize}

\paragraph{Validation of the per-dish scale.}
$s_f$ is only valid if \texttt{ROIdef}'s radius is in the same pixel frame as the
fish positions. We verified this by overlaying each fish's trajectory with a
circle of its ROI radius: the explored extent (99th-percentile radial distance
from the fish's median position) is a consistent $0.96\pm0.01$ of the ROI radius
for all 15 fish, confirming the radius is in the tracking-pixel frame (each fish
explores $\sim$96\% of its dish). Note this per-dish scaling is \emph{not} in the
original repo, which keeps fish positions in pixels throughout; it is introduced
here so that bout distances are in millimetres and comparable across fish.

\plotph{Per-fish trajectory extent vs.\ ROIdef radius -- 15-panel overlay
        (notebook Sect.~3 validation).}


\section{Bout discretization scheme}

\subsection{Overview and origin}
Locomotion is reduced from the continuous 30~fps track to a sequence of discrete
swim \emph{bouts}, each summarised by three kinematic scalars: reorientation
angle $\delta\theta$, inter-bout interval $\delta t$, and travelled distance $d$.
The detection method follows Le Goc et al.\ 2021 (BMC Biology,
\texttt{10.1186/s12915-021-01126-w}), with all time-based thresholds reconverted
from their 25~fps to our 30~fps.

\paragraph{Smoothing for detection, raw positions for measurement.}
The trajectory is Gaussian-smoothed \emph{only to detect bout timing}: the
integer-pixel jitter ($\approx0.4$~mm) otherwise spawns spurious speed peaks.
Everything \emph{measured} from a bout -- heading, reorientation, distance,
onset position -- is read from the \textbf{raw (unsmoothed)} positions, so the
smoothing kernel never biases the reported direction.

\paragraph{Two anchors per bout.}
Each detected bout carries \emph{two} reference frames:
\begin{itemize}
  \item the \textbf{onset} $o_i$ -- the speed \emph{trough} just before the bout
        (the rest point from which the fish accelerates). It marks the physical
        bout start and is the anchor for \emph{timing}, \emph{position},
        \emph{stimulus context}, and -- computed \emph{onset-to-onset}, see
        Step~7 -- for \emph{heading}, \emph{reorientation} and \emph{distance};
  \item the \textbf{peak} $p_i$ -- the speed \emph{maximum}. It is kept only as
        an \emph{independent cross-check}: the tracked body orientation at the
        peak (flip-corrected) gives a second estimate $\delta\theta^{\text{ori}}$
        of the reorientation, used to validate the displacement-based one.
\end{itemize}
Heading is taken onset-to-onset (net displacement between two rest points)
rather than from an instantaneous single-frame direction, because on raw
integer-pixel data a single rest-frame direction is dominated by quantisation
noise, whereas the net leg between successive onsets is a robust travel
direction.

\subsection{Parameters and thresholds}
\begin{table}[h]\centering\small
\begin{tabular}{@{}lll@{}}
\toprule
\textbf{Parameter} & \textbf{Value (30~fps)} & \textbf{Origin}\\
\midrule
Frame rate $f_s$                  & 30~fps                  & acquisition\\
Position smoothing $\sigma_{\text{pos}}$ & 2 frames (Gaussian, detection only) & noise control\\
Speed smoothing window            & 3 frames (rolling mean) & noise control\\
Speed-threshold factor $k$        & 1.0 (tunable)           & this work\\
Bout onset                        & speed trough before peak (default) & this work\\
Heading / distance                & onset-to-onset, raw positions & this work\\
Onset ``lead'' option             & 2 frames ($\approx$80~ms) before peak & Le Goc\\
Refractory (min.\ peak spacing)   & 12 frames (0.4~s)       & Le Goc\\
Displacement window $\pm w$       & 15 frames ($\pm$0.5~s)  & Le Goc\\
Displacement gate $[d_{\min},d_{\max}]$ & $[0.3,\,18]$~mm   & Le Goc\\
Forward threshold $\delta\theta_0$ & $10^\circ$             & Le Goc / HMM\\
Border shell                      & 5~mm from dish edge     & Le Goc ROI\\
\bottomrule
\end{tabular}
\end{table}

\subsection{Step-by-step procedure}
The procedure is applied independently to each fish. Below, $x,y$ denote the
\emph{raw} scaled positions and $\tilde x,\tilde y$ the smoothed ones (used only
in Steps~4--6 for detection).

\paragraph{Step 1 -- Scale to millimetres.}
Convert positions with the per-dish factor $s_f$:
$\;x_{\text{mm}} = s_f\, x_{\text{px}},\; y_{\text{mm}} = s_f\, y_{\text{px}}$.

\paragraph{Step 2 -- Smooth the trajectory (detection only).}
Each coordinate time series is convolved with a one-dimensional Gaussian kernel
$G_\sigma$ of standard deviation $\sigma_{\text{pos}}=2$~frames:
\[
\tilde x_t = \sum_{k=-K}^{K} G_\sigma(k)\, x_{t-k},
\qquad
G_\sigma(k) = \frac{1}{Z}\,\exp\!\Big(-\frac{k^2}{2\sigma_{\text{pos}}^2}\Big),
\]
and likewise for $\tilde y$, where $k$ indexes frame offsets, $Z=\sum_k
\exp(-k^2/2\sigma_{\text{pos}}^2)$ normalises the weights to unit sum, and the
kernel is truncated at $K=\lceil 4\sigma_{\text{pos}}\rceil$ (the default support
of \texttt{scipy.ndimage.gaussian\_filter1d}, which we use). This restores a
continuous, sub-pixel path so that the speed peaks used for detection are not
corrupted by integer-pixel jitter. Because $\sigma_{\text{pos}}=2$~frames
($\approx67$~ms) is short relative to the inter-bout interval ($\sim$1~s), it
removes jitter without blurring distinct bouts together. \emph{The smoothed
series is used only in Steps~4--6; heading, reorientation, distance and position
(Step~7 onward) are read from the raw $x,y$.}

\paragraph{Step 3 -- Segment continuous pieces.}
Break the per-fish stream into continuous segments \textbf{only where tracking is
lost} (NaN gaps). Trajectories are \emph{not} cut at stimulus boundaries, so the
bout in which the fish first reacts to a newly-appeared stimulus keeps its
reorientation; each bout is later tagged by the condition active at its onset.

\paragraph{Step 4 -- Instantaneous speed.}
Per segment, form the frame-to-frame speed from the \emph{smoothed} path and
lightly smooth it (3-frame rolling mean):
\[
v(t) \;=\; f_s\,\sqrt{\big(\tilde x_{t}-\tilde x_{t-1}\big)^2
                      +\big(\tilde y_{t}-\tilde y_{t-1}\big)^2}\quad[\text{mm\,s}^{-1}].
\]

\paragraph{Step 5 -- Detect bouts (peaks) and place onsets.}
A bout is a local maximum of $v$ exceeding a data-driven threshold, with peaks
forced at least one refractory period apart:
\[
v(p_i) > \mu_v + k\,\sigma_v ,\qquad
\mu_v=\overline{v},\;\; \sigma_v=\mathrm{std}(v),
\]
consecutive peaks $\ge 12$ frames apart. The bout \emph{onset} $o_i$ is then
placed at the local speed minimum preceding $p_i$ (the trough; ``walk back from
$p_i$ until the speed stops decreasing''). A Le~Goc-literal alternative places it
$o_i = p_i - 2$ frames ($\approx$80~ms before the peak).

\paragraph{Step 6 -- Displacement gate.}
Reject bouts whose $\pm 0.5$~s net displacement around the onset is too small
(wobble) or too large (tracking jump):
\[
d_{\min} \;\le\;
\big\lVert \mathbf{r}(o_i+w)-\mathbf{r}(o_i-w)\big\rVert \;\le\; d_{\max},
\qquad d_{\min}=0.3,\; d_{\max}=18~\text{mm}.
\]

\plotph{Speed trace $v(t)$ over one stimulus block, with onsets (troughs) and
        speed peaks marked (notebook \texttt{segment\_bouts.ipynb}, Sect.~5).}

\paragraph{Step 7 -- Per-bout kinematics (onset-to-onset, raw positions).}
With ordered onsets $o_i$, the \emph{travel direction} (heading) of bout $i$ is
the direction of the net \emph{onset-to-onset} displacement, measured on the raw
positions:
\[
\theta_i = \operatorname{atan2}\!\big(y_{o_{i+1}}-y_{o_i},\;
                                      x_{o_{i+1}}-x_{o_i}\big).
\]
The \emph{reorientation angle} is the wrapped change of heading:
\[
\delta\theta_i = \operatorname{wrap}\!\big(\theta_i-\theta_{i-1}\big),\qquad
\operatorname{wrap}(a)=\big((a+\pi)\bmod 2\pi\big)-\pi \in(-\pi,\pi].
\]
The \emph{inter-bout interval} and \emph{travelled distance} are the matching
onset-to-onset quantities:
\[
\delta t_i = \frac{o_{i+1}-o_i}{f_s},\qquad
d_i = \big\lVert \mathbf{r}_{o_{i+1}}-\mathbf{r}_{o_i}\big\rVert .
\]
As an independent check we also store $\delta\theta^{\text{ori}}$, computed from
the tracked body orientation $f\_ori$ \emph{at the peak} after flip-correcting it
to the travel direction ($\theta^{\text{ori}} \mapsto \theta^{\text{ori}}+\pi$
when $\cos(\theta^{\text{ori}}-\theta)<0$).

\plotph{2D trajectory coloured by detected bouts, with onset-to-onset heading
        arrows (notebook Sect.~6).}

\paragraph{Step 8 -- Stimulus context.}
Tag each bout with the \texttt{stimulus} label at its \emph{onset}, parsed into
$n_{\text{left}}$, $n_{\text{right}}$, the side, and dot speeds. The dot-group
angle relative to the fish facing direction (evaluated at the onset) is
\[
\phi_i = \operatorname{wrap}\!\Big(
\operatorname{atan2}\big(\bar y_{\text{dots}}-y_{o_i},\,
                         \bar x_{\text{dots}}-x_{o_i}\big)-f\_ori_{o_i}\Big),
\]
with $(\bar x_{\text{dots}},\bar y_{\text{dots}})$ the centroid of the active
dots and the sign convention $\phi_i>0 \Rightarrow$ dots on the fish's left.
$\phi_i$ (\texttt{stim\_angle\_rad}) is the driving input of the stimulus-aware
model in Sect.~4.

\paragraph{Step 9 -- Bout-type label.}
Classify by the $10^\circ$ threshold:
\[
\text{type}_i=
\begin{cases}
\text{forward}, & |\delta\theta_i| < 10^\circ,\\
\text{left},    & \delta\theta_i \ge 10^\circ,\\
\text{right},   & \delta\theta_i \le -10^\circ.
\end{cases}
\]

\paragraph{Step 10 -- Border flag (non-destructive).}
Mirroring Le Goc's 5~mm ROI, flag (do not drop) bouts in the outer shell. The
dish centre and radius are estimated empirically per fish (median position;
95th-percentile radial distance $R_f$), using the onset position:
\[
\texttt{near\_wall}_i = \Big[\;\big\lVert \mathbf{r}_{o_i}-\mathbf{c}_f\big\rVert
> R_f - 5~\text{mm}\;\Big].
\]

\subsection{Output dataframe}
One row per detected bout, with columns:
\begin{table}[h]\centering\small
\begin{tabular}{@{}ll@{}}
\toprule
\textbf{Column} & \textbf{Meaning}\\
\midrule
\texttt{fish\_id, segment\_id, bout\_index} & fish, continuous segment, index within segment\\
\texttt{frame\_onset, t\_onset\_s}          & bout onset (trough; frame and seconds)\\
\texttt{dtheta\_rad, dtheta\_deg}           & reorientation $\delta\theta$ (onset-to-onset, raw)\\
\texttt{dtheta\_ori\_deg}                   & orientation-based cross-check (at peak)\\
\texttt{dt\_s}                              & inter-bout interval (onset-to-onset)\\
\texttt{distance\_mm}                       & travelled distance (onset-to-onset, raw)\\
\texttt{heading\_rad, f\_ori\_rad}          & travel direction; raw body orientation\\
\texttt{x\_mm, y\_mm}                       & onset (trough) position, raw (mm)\\
\texttt{stimulus, stim\_angle\_rad}         & condition at onset; dot-group bearing $\phi$\\
\texttt{n\_dots\_left, n\_dots\_right}      & group size per side\\
\texttt{speed\_left, speed\_right, moving}  & dot speeds; moving flag\\
\texttt{side}                               & left / right / both / none\\
\texttt{bout\_type, near\_wall}             & F/L/R label; border flag\\
\bottomrule
\end{tabular}
\end{table}

\subsection{Deviations from Le Goc et al.}
\begin{itemize}
  \item \textbf{Smoothing for detection only}: bout \emph{timing} is found on the
        Gaussian-smoothed track, but heading, $\delta\theta$, distance and
        position are measured on the \emph{raw} positions, onset-to-onset. This
        keeps the reported direction free of any smoothing bias.
  \item \textbf{Heading from displacement}, not body axis: the tracked
        \texttt{f\_ori} here suffers $180^\circ$ head/tail flips, so the
        onset-to-onset travel direction is more robust; \texttt{f\_ori} is kept
        only as the cross-check $\delta\theta^{\text{ori}}$.
  \item \textbf{Split only at tracking gaps} (not at stimulus blocks), to keep
        the onset reorientation toward newly-appeared stimuli.
  \item \textbf{Per-dish pixel$\to$mm scale} (from \texttt{ROIdef}), not a single
        global factor.
  \item \textbf{No turn-displacement ($1/1.6$) correction}: it targets
        temperature comparisons absent here, the $1.6$ factor is
        setup-specific, and it is dominated by the resolution-driven distance
        inflation.
\end{itemize}

\subsection{Validation}
The segmented kinematics are compared against the freely-swimming HMM2023
reference data (Le Goc data, 26~$^\circ$C): bout rate, distance and $\delta\theta$
distributions. Our bouts remain coarser than the reference (lower rate, longer
legs, broader $\delta\theta$: median $|\delta\theta|\approx39^\circ$ and forward
fraction $\approx0.11$, vs.\ $\approx10^\circ$ and $0.50$ in the reference),
consistent with adjacent bouts merging at the low integer-pixel tracking
resolution and with the onset-to-onset (rather than peak-to-peak) heading, which
we retain because it is the direction the fish actually travels until the next
bout. The signed $\delta\theta$ is unbiased (no left/right handedness artefact).

\plotph{Reorientation $\delta\theta$ histogram: displacement-derived vs.\
        flip-corrected orientation (notebook Sect.~9).}
\plotph{Inter-bout interval and bout-distance distributions (notebook Sect.~9).}
\plotph{Reorientation by stimulus side -- left vs.\ right dots (notebook
        Sect.~10).}
\plotph{Overlay vs.\ HMM2023 paper data ($\delta\theta$, $\delta t$, distance)
        (notebook \texttt{compare\_to\_paper.ipynb}).}


\section{Behavioural-state models}
\label{sec:hmm}

\subsection{Rationale and state alphabet}
Following Dommanget-Kott et al.\ 2026 (\emph{Linking brain and behavior states in
zebrafish larvae locomotion using Hidden Markov Models}), we model the sequence
of reorientation angles $\{\delta\theta_n\}$ produced by the discretization as
emissions of a small number of latent locomotor states. We fit two nested
models: (i) a \emph{Markov chain} on hard-thresholded states (baseline), and
(ii) a 3-state \emph{hidden Markov model} (HMM) with continuous emissions. Both
use three states -- one \emph{forward} and two \emph{turning} states -- but the
HMM infers them probabilistically rather than by a hard cut.

Throughout, the two turning states are written \textbf{sign-neutrally} as
\emph{turn$-$} (emitting $\delta\theta<0$) and \emph{turn$+$} (emitting
$\delta\theta>0$); ``left/right'' labels are avoided so that the sign of a turn is
unambiguous, which matters for the stimulus coupling of Sect.~\ref{sec:iohmm}.

\subsection{Markov chain baseline (thresholded)}
Each bout is hard-classified with the forward threshold $\delta\theta_0=10^\circ$
into a state $s_n\in\{F,L,R\}$,
\[
s_n=
\begin{cases}
F, & |\delta\theta_n|<\delta\theta_0,\\
L, & \delta\theta_n \ge \delta\theta_0 \quad(\text{turn}+),\\
R, & \delta\theta_n \le -\delta\theta_0 \quad(\text{turn}-),
\end{cases}
\]
reproducing the paper's \texttt{utils/MarkovChains.py}. The transition matrix is
estimated by counting and row-normalising,
\[
P_{ij}=\frac{N_{ij}}{\sum_{k}N_{ik}},
\qquad
N_{ij}=\#\{\,n:\ s_n=i,\ s_{n+1}=j\,\},
\]
and its stationary distribution $\boldsymbol\pi$ (the a-priori state
occupancies) solves
\[
\boldsymbol\pi P = \boldsymbol\pi,\qquad \textstyle\sum_i \pi_i = 1 .
\]
On our data the forward occupancy is low ($\pi_F\approx0.11$) and the two turn
states are near-symmetric ($\pi_L\approx\pi_R\approx0.44$), consistent with the
broad $\delta\theta$ distribution.

\plotph{MC transition matrix (heat map) and $\delta\theta$ histogram with the
        $\pm10^\circ$ thresholds (notebook \texttt{fit\_hmm.ipynb}, Sect.~2).}

\subsection{Hidden Markov model (\texttt{G3\_Sym})}
\paragraph{Emissions.}
The reorientation of state $j$ is drawn from a continuous density
$b_j(\delta\theta)$ (angles in degrees):
\[
\begin{aligned}
b_F(\delta\theta)      &= \mathcal N(\delta\theta;\,0,\sigma)
   = \tfrac{1}{\sqrt{2\pi\sigma^2}}\,e^{-\delta\theta^2/2\sigma^2},\\
b_{\text{turn}+}(\delta\theta) &=
   \Gamma(\delta\theta;\,\alpha,\theta)\,\mathbf 1[\delta\theta>0],\qquad
b_{\text{turn}-}(\delta\theta) =
   \Gamma(-\delta\theta;\,\alpha,\theta)\,\mathbf 1[\delta\theta<0],
\end{aligned}
\]
with the Gamma density
$\Gamma(x;\alpha,\theta)=x^{\alpha-1}e^{-x/\theta}/(\Gamma(\alpha)\,\theta^{\alpha})$.
Forward bouts are a centred Normal (small, symmetric wobble); turns are a
\emph{single, shared} Gamma over the turn magnitude, reflected onto each sign.
This is exactly the emission set stored in the paper's fitted models
($\sigma_{\text{forw}}$ plus $\texttt{turn}=[\alpha,\theta]$), with the
constraints $\alpha\ge1$, $\theta>0$.

\paragraph{Left--right symmetry.}
Transitions and the initial distribution are tied to be left/right symmetric:
$P(F\!\to\!L)=P(F\!\to\!R)$, $P(L\!\to\!L)=P(R\!\to\!R)$,
$P(L\!\to\!R)=P(R\!\to\!L)$, $P(L\!\to\!F)=P(R\!\to\!F)$, and the turn states
share $(\alpha,\theta)$. This encodes the absence of an intrinsic turning bias.

\paragraph{Likelihood.}
For a trajectory $\delta\theta_{1:N}$ with hidden states $s_{1:N}$, transition
matrix $A$ and initial distribution $\boldsymbol\pi$,
\[
\mathcal L
=\sum_{s_{1:N}} \pi_{s_1}\,b_{s_1}(\delta\theta_1)\;
\prod_{n=2}^{N} A_{s_{n-1}s_n}\,b_{s_n}(\delta\theta_n).
\]

\paragraph{Optimiser: Baum-Welch (EM).}
The parameters $\Theta=(\boldsymbol\pi,A,\sigma,\alpha,\theta)$ are fit by
Baum-Welch, the Expectation-Maximisation algorithm for HMMs. This is the
optimiser used by the paper (a generic Baum-Welch specialised to this emission
family); we reimplement it directly. Define the forward/backward variables
\[
\alpha_n(i)=P(\delta\theta_{1:n},s_n=i),\qquad
\beta_n(i)=P(\delta\theta_{n+1:N}\mid s_n=i),
\]
computed by the usual recursions (in practice in the scaled form of Rabiner to
avoid underflow). The \textbf{E-step} posteriors are
\[
\gamma_n(i)=\frac{\alpha_n(i)\beta_n(i)}{\mathcal L},\qquad
\xi_n(i,j)=\frac{\alpha_n(i)\,A_{ij}\,b_j(\delta\theta_{n+1})\,\beta_{n+1}(j)}{\mathcal L}.
\]
The \textbf{M-step} re-estimates each parameter in closed form (then applies the
symmetry tying):
\[
A_{ij}=\frac{\sum_n \xi_n(i,j)}{\sum_n \gamma_n(i)},\qquad
\pi_i=\gamma_1(i),\qquad
\sigma^2=\frac{\sum_n \gamma_n(F)\,\delta\theta_n^2}{\sum_n \gamma_n(F)} .
\]
The shared turn Gamma is fit by a \emph{weighted} maximum-likelihood on the turn
magnitudes $m_n=|\delta\theta_n|$ with weights
$w_n=\gamma_n(\text{turn}+)\,\mathbf 1[\delta\theta_n>0]
     +\gamma_n(\text{turn}-)\,\mathbf 1[\delta\theta_n<0]$. Writing
$W=\sum_n w_n$, $\bar m=\tfrac1W\sum_n w_n m_n$,
$\overline{\log m}=\tfrac1W\sum_n w_n\log m_n$ and
$s=\log\bar m-\overline{\log m}$, the shape is obtained from Minka's
approximation refined by Newton,
\[
\alpha \approx \frac{3-s+\sqrt{(s-3)^2+24\,s}}{12\,s},\qquad
\theta=\frac{\bar m}{\alpha},\qquad \alpha\leftarrow\max(\alpha,1),
\]
and the scale follows as $\theta=\bar m/\alpha$. EM is iterated until the
log-likelihood gain falls below $10^{-6}$ (max 500 iterations), from several
random left--right-symmetric initialisations; the best run is kept.

\paragraph{Fitted values.}
On \texttt{GsizeMotionSingle} the fit is stable across restarts and gives
$\sigma\approx6.5^\circ$, $\alpha\approx2.3$, $\theta\approx21^\circ$ (mean turn
$\alpha\theta\approx49^\circ$), with a sticky, symmetric transition matrix. The
forward emission is much broader than the paper's clean data
($\sigma\approx3.5^\circ$ at 26~$^\circ$C), as expected from the coarser
tracking.

\plotph{Fitted emission densities (forward Normal + signed turn Gamma) overlaid
        on the $\delta\theta$ histogram (notebook \texttt{fit\_hmm.ipynb},
        Sect.~4).}
\plotph{Viterbi state labelling of a real trajectory, MC threshold vs.\ HMM
        (notebook \texttt{fit\_hmm.ipynb}, Sect.~5).}

\subsection{Trajectory generation}
\label{sec:gen}
To generate synthetic paths (paper Sect.~2.7) we (i) sample a state sequence
$s_{1:N}$ from $(\boldsymbol\pi,A)$; (ii) sample a reorientation
$\delta\theta_n\sim b_{s_n}$ from the corresponding emission; and (iii) resample
the inter-bout interval $\delta t_n$ and distance $d_n$ \emph{independently} from
the empirical pools of the data (ignoring any bout-type dependence, as the paper
does). The 2D path follows by cumulative summation of heading and steps:
\[
\theta_n=\theta_0+\sum_{k\le n}\delta\theta_k,\qquad
x_n=\sum_{k\le n} d_k\cos\theta_k,\qquad
y_n=\sum_{k\le n} d_k\sin\theta_k .
\]
The generated angle distribution matches the data (median $|\delta\theta|$
$38.8^\circ$ vs.\ $39.0^\circ$).

\plotph{HMM-generated synthetic trajectories and real-vs-generated
        $\delta\theta/\delta t/d$ distributions (notebook
        \texttt{fit\_hmm.ipynb}, Sect.~6--7).}


\section{Stimulus-driven (input-driven) HMM}
\label{sec:iohmm}

\subsection{Rationale}
The spontaneous HMM ignores the stimulus. To capture the observed
\emph{following} of the moving dots (Perrino thesis: fish orient toward and
approach conspecific-like motion), we extend it to an \textbf{input-driven HMM}
(IO-HMM / GLM-HMM), in which the transition and the forward emission depend, at
each bout, on the egocentric bearing to the active dots. We first verify that the
effect is present in the raw data: among moving-stimulus bouts that are turns,
$62\%$ turn \emph{toward} the dots (vs.\ $50\%$ chance).

\subsection{Driving input}
The bearing $\phi_n$ (\texttt{stim\_angle\_rad}, Step~8) is turned into a bounded
following signal
\[
b_n=\sin(\phi_n)\cdot\mathbf 1[\text{active moving stimulus at bout }n]
    \in[-1,1],
\]
($b_n=0$ when no moving dot is present). $\sin\phi$ gives the sign of the turn
needed to face the dots and vanishes when they are straight ahead or behind.
Crucially, because $\delta\theta$ and $\phi$ share the same world frame, turning
\emph{toward} the dots always means $\operatorname{sign}(\delta\theta)=\operatorname{sign}(\phi)$;
so with sign-neutral states, $b_n>0$ (dots on the left) must favour turn$+$.
During generation $b_n$ is \emph{recomputed each bout} from the fish's current
pose and the real dot positions, so the synthetic fish reacts as it swims.

\subsection{Model}
Let $c=(c_F,c_{-},c_{+})=(0,-1,+1)$ index the sign push of each state. The
\textbf{input-driven transition} from state $i$ is a multinomial logit whose turn
columns are tilted by $b_n$,
\[
A^{(n)}_{ij}=\frac{\exp\!\big(\log T^{0}_{ij}+\kappa\,c_j\,b_n\big)}
                  {\sum_{k}\exp\!\big(\log T^{0}_{ik}+\kappa\,c_k\,b_n\big)},
\]
with $T^0$ the (symmetric) baseline transition and $\kappa\ge0$ the learned
\emph{following gain}. When $b_n=0$ this reduces to $A^{(n)}=T^0$, recovering the
spontaneous HMM; when $b_n>0$ it multiplies $P(\text{turn}+)$ by $e^{+\kappa b_n}$
and $P(\text{turn}-)$ by $e^{-\kappa b_n}$, biasing the turn toward the dots. The
\textbf{forward emission} additionally drifts toward the stimulus,
\[
b_F(\delta\theta;b_n)=\mathcal N(\delta\theta;\,g\,b_n,\ \sigma),
\]
with gain $g$; the turn emissions are unchanged. The full parameter set is
$\Theta=(T^0,\kappa,\,\sigma,g,\,\alpha,\theta,\boldsymbol\pi)$.

\subsection{Optimiser: input-driven Baum-Welch}
The EM skeleton is unchanged; only the transition/emission are now functions of
$b_n$, so the forward--backward uses a \textbf{time-varying} transition
$A^{(n)}=A(b_n)$ (and $b_F(\cdot;b_n)$). The E-step recursions and the posteriors
$\gamma_n(i),\ \xi_n(i,j)$ keep the same form as above, with $A_{ij}$ replaced by
$A^{(n-1)}_{ij}$ where appropriate.

The M-step splits into a closed-form part and a small numerical part:
\begin{itemize}
  \item \textbf{Turn Gamma} $(\alpha,\theta)$: the same weighted MLE as the
        spontaneous model -- independent of $b$.
  \item \textbf{Forward drift and variance}: weighted least-squares through the
        origin and weighted variance,
        \[
        g=\frac{\sum_n \gamma_n(F)\,b_n\,\delta\theta_n}
               {\sum_n \gamma_n(F)\,b_n^{2}},\qquad
        \sigma^2=\frac{\sum_n \gamma_n(F)\,(\delta\theta_n-g\,b_n)^2}
                      {\sum_n \gamma_n(F)} .
        \]
  \item \textbf{Baseline transition and following gain} $(T^0,\kappa)$: no closed
        form. They maximise the expected complete-data transition
        log-likelihood
        \[
        Q(T^0,\kappa)=\sum_n\sum_{i,j}\xi_n(i,j)\,
        \log A^{(n)}_{ij}(T^0,\kappa),
        \]
        i.e.\ a small multinomial-logit (GLM-HMM) fit, solved with L-BFGS. To
        keep it cheap, the responsibilities $\xi_n$ are \emph{binned by the
        scalar input} $b_n$ ($\sim$40 bins on $[-1,1]$), so the objective is a
        sum over bins rather than over all $\sim7\times10^{4}$ transitions.
\end{itemize}
This is the standard GLM-HMM / IO-HMM M-step. EM converges under the same
$10^{-6}$ criterion and is run from several restarts.

\subsection{Fitted result}
The fit is stable across restarts, with a slightly higher log-likelihood than the
spontaneous model (as expected from the extra parameters), and recovers the same
spontaneous parameters ($\sigma\approx6.5^\circ,\ \alpha\approx2.3,\
\theta\approx21^\circ$, symmetric $T^0$) plus a \textbf{positive following gain}
$\kappa\approx0.40>0$: the fish turns toward the dots. At full salience
($b=1$, dots at $90^\circ$) this makes a turn-toward $\approx2.2\times$ more
likely than a turn-away from a turn state. The forward drift $g$ is small, so the
following acts mainly through the transition.

\plotph{Following bias: $P(\text{turn}\pm\mid\text{turn state})$ vs.\ input
        $b=\sin\phi$ (the $\kappa$ effect), and emissions vs.\ data (notebook
        \texttt{fit\_hmm\_stim.ipynb}, Sect.~3).}

\subsection{Stimulus-aware generation and comparison video}
Synthetic paths are generated as in Sect.~\ref{sec:gen}, but each bout's state
and forward mean use $A^{(n)}$ and $b_F(\cdot;b_n)$ with $b_n$ recomputed from the
synthetic pose and the real dot centroid, so the fish steers toward the dots. In
a controlled test with a fixed virtual stimulus, the fitted model spends $38\%$
of its time within $30$~mm of the stimulus vs.\ $16\%$ at $\kappa=0$ -- following
more than doubles time-near-stimulus.

For a like-for-like video, the bout-level output is rendered back to a per-frame
($30$~fps) trajectory: the fish holds its pose during the inter-bout interval and
executes the sampled reorientation and step at the bout onset (``wait, then
dart''); distances are converted to pixels with the per-dish scale
$s_f^{-1}$. Each synthetic segment is seeded from the real fish for the first
few seconds and carries the \emph{real} dot positions, so a 2$\times$N video
shows the real fish (top) and the stimulus-driven HMM (bottom) under exactly the
same stimuli.

\plotph{Controlled following demo: paths for $\kappa{=}0$ vs.\ fitted $\kappa$
        with a fixed stimulus (notebook \texttt{fit\_hmm\_stim.ipynb}, Sect.~4).}
\plotph{Comparison video frames: real (top) vs.\ stimulus-driven HMM (bottom)
        over the same dots (\texttt{runs/hmm\_stim/comparison.mp4}).}

\end{document}
